\documentclass[12pt]{article}
\usepackage[utf8]{inputenc}
\usepackage[T1]{fontenc}
\usepackage[margin=1in]{geometry}
\usepackage{setspace}
\usepackage{xcolor}
\usepackage[colorlinks=true, urlcolor=blue!60!black, linkcolor=blue!60!black]{hyperref}
\setlength{\parindent}{2em}
\setlength{\parskip}{1.25em}
\renewcommand{\baselinestretch}{1.0}
\begin{document}
\begin{center}
{\LARGE \textit{Cuyler v.\ Adams}, 449 U.S.\ 433 (1981)\par}
\vspace{0.8em}
{\large Case Brief --- Four-Page Edition\par}
\end{center}
\vspace{1em}

\section{Memory Jogger}

When Congress consents, even in advance, to an interstate compact on a subject appropriate for federal legislation, the compact becomes federal law. The Court held 6--3 that the Interstate Agreement on Detainers is federal law and that its Article IV(d) preserves a prisoner's extradition-hearing rights (\href{https://supreme.justia.com/cases/federal/us/449/433/}{Justia}).

\section{Detailed Case Facts}

In April 1976 John Adams was convicted of robbery in Pennsylvania and sentenced to 30 years at the State Correctional Institution at Graterford. The Camden County, New Jersey, prosecutor lodged a detainer against him and in May 1977 filed a ``Request for Temporary Custody'' under Article IV of the Interstate Agreement on Detainers (IAD) to bring him to New Jersey for trial on armed-robbery charges. The IAD is a compact among 48 States, the District of Columbia, and the United States. Both Pennsylvania and New Jersey had also adopted the Uniform Criminal Extradition Act (UCEA), whose \S~10 guarantees a pre-transfer judicial hearing.

\section{Procedural History}

Adams filed a \textit{pro se} class action under 42 U.S.C.\ \S\S~1981 and 1983, alleging that Pennsylvania officials denied him the UCEA hearing and failed to tell him of his right under Article IV(a) to petition the Governor to disapprove the transfer. The district court dismissed. Adams was transferred, convicted in New Jersey, and returned. Meanwhile a Pennsylvania appellate court held in \textit{Commonwealth ex rel.\ Coleman v.\ Cuyler} that Article IV transferees have no right to a pre-transfer hearing. The Third Circuit nonetheless vacated the dismissal, holding that the IAD, as a congressionally sanctioned compact, is federal law and that Article IV(d) incorporates the UCEA's safeguards. The Supreme Court granted certiorari to resolve a split.

\section{Judicial Votes}

\textbf{6--3}, affirmed (January 21, 1981).
\begin{itemize}
\item \textbf{Majority:} Justice Brennan, joined by Justices White, Marshall, Blackmun, Powell, and Stevens.
\item \textbf{Concurrences:} None.
\item \textbf{Dissent:} Justice Rehnquist, joined by Chief Justice Burger and Justice Stewart.
\end{itemize}

\section{Holding}

Under the ``law of the Union'' doctrine, where Congress authorizes States to enter a cooperative agreement on a subject appropriate for congressional legislation, its consent transforms the agreement into federal law, and its construction is a federal question. Congress consented in advance to the IAD in the Crime Control Consent Act of 1934, on a proper federal subject, so the IAD is federal law. Article IV(d) preserves a prisoner's UCEA safeguards, including a pre-transfer hearing, and Adams stated a \S~1983 claim.

\section{Analysis of Opinions}

\textbf{Majority (Brennan): federal law.} From \textit{Wheeling Bridge} through \textit{Petty v.\ Tennessee-Missouri Bridge Comm'n}, congressional consent has made compacts federal law. \textit{Virginia v.\ Tennessee} and \textit{U.S.\ Steel v.\ Multistate Tax Comm'n} hold that consent is \textit{required} only for agreements that increase state power at federal expense, but when Congress \textit{does} consent on a proper subject, the agreement is federalized. Consent may come in advance, as in the 1934 Act, reaffirmed when the United States joined the IAD in 1970. Detainers fall within the Commerce and Extradition Clauses.

\textbf{Majority (Brennan): Article IV(d).} Unbound by \textit{Coleman}, the Court contrasted Article III, where a prisoner's own request is deemed a ``waiver of extradition,'' with Article IV's prosecutor-initiated transfers. Article IV(d) preserves ``any right which he may have to contest the legality of his delivery,'' except objection based on the Governor's lack of consent. Reading it to protect only the petition right would make it superfluous; withholding one UCEA defense implies the others survive; and ``any right'' signals several. The drafting history confirmed the remedial purpose. A federal statutory right supported a \S~1983 claim under \textit{Maine v.\ Thiboutot}.

\textbf{Rehnquist, joined by Burger and Stewart (dissenting).} The majority's ``remarkable feat of judicial alchemy'' turned a Pennsylvania statute into federal law. Because the IAD does not threaten federal supremacy, consent was unnecessary, and gratuitous advance consent cannot federalize state law; consent is not legislation (\textit{Coyle v.\ Smith}; \textit{Prudential v.\ Benjamin}). The ruling threatened to become a ``judicial Midas'' federalizing uniform state criminal statutes. He would have deferred to \textit{Coleman}.

\textbf{Later developments.} \textit{Texas v.\ New Mexico}, \textit{Alabama v.\ North Carolina}, \textit{Tarrant Regional Water District}, and \textit{Kansas v.\ Nebraska} reaffirmed that a consented compact is federal law; \textit{Carchman v.\ Nash} and \textit{Reed v.\ Farley} built on \textit{Cuyler} in IAD cases.

\section{Examples: Future Applications}

\begin{enumerate}
\item \textbf{Same-Side: Disaster Compact.} A state court's narrow reading of the congressionally consented Emergency Management Assistance Compact's immunity clause is not binding; the dispute is a federal question.
\item \textbf{Same-Side: Bridge Commission.} One State cannot impose its bidding rules on a consented bistate bridge commission; the compact's meaning is federal.
\item \textbf{Opposite-Side: Cosmetology Compact.} An unconsented licensing compact that does not threaten federal supremacy is valid without consent and remains state law, creating no \S~1983 right.
\item \textbf{Opposite-Side: Technical IAD Violation.} An unobjected-to, harmless breach of the 120-day trial deadline does not support federal habeas relief (\textit{Reed v.\ Farley}).
\item \textbf{Fence-Sitter: Uniform Witness Act.} Whether the 1934 advance consent federalizes the Uniform Act to Secure the Attendance of Witnesses is unresolved; Rehnquist raised it as an example of the majority's reach.
\end{enumerate}

\section{Critique}

\textbf{Consent versus enactment.} Rehnquist's critique retains bite: federal statutes require bicameralism and presentment of specific text, and treating a 1934 blanket consent as transforming a 1956 state-drafted uniform law into federal law blurs the line between congressional permission and lawmaking.

\textbf{Scope of advance consent.} Footnote 10 declined to decide whether other reciprocal criminal statutes are likewise federalized, leaving lower courts to separate formal compacts from ordinary reciprocal legislation.

\section{Key Quotations}

\begin{quote}
``But where Congress has authorized the States to enter into a cooperative agreement, and where the subject matter of that agreement is an appropriate subject for congressional legislation, the consent of Congress transforms the States' agreement into federal law under the Compact Clause.''
\end{quote}
\noindent --- Justice Brennan, for the Court, 449 U.S.\ at 440.

\begin{quote}
``Because congressional consent transforms an interstate compact within this Clause into a law of the United States, we have held that the construction of an interstate agreement sanctioned by Congress under the Compact Clause presents a federal question.''
\end{quote}
\noindent --- Justice Brennan, for the Court, 449 U.S.\ at 438.

\begin{quote}
``By vesting in Congress the power to grant or withhold consent, or to condition consent on the States' compliance with specified conditions, the Framers sought to ensure that Congress would maintain ultimate supervisory power over cooperative state action that might otherwise interfere with the full and free exercise of federal authority.''
\end{quote}
\noindent --- Justice Brennan, for the Court, 449 U.S.\ at 439--440.

\end{document}
